\documentclass[10pt, letterpaper]{article}

% Packages:
\usepackage[
    ignoreheadfoot,
    top=0.4 cm,
    bottom=0.4 cm,
    left=0.8 cm,
    right=0.8 cm,
    footskip=0.2 cm,
]{geometry}
\usepackage{titlesec}
\usepackage{tabularx}
\usepackage{array}
\usepackage[dvipsnames]{xcolor}
\definecolor{primaryColor}{RGB}{0, 0, 0}
\usepackage{enumitem}
\usepackage{amsmath}
\usepackage[
    pdftitle={Tanmay Bhuskute's CV},
    pdfauthor={Tanmay Bhuskute},
    pdfcreator={LaTeX with RenderCV},
    colorlinks=true,
    urlcolor=primaryColor
]{hyperref}
\usepackage[pscoord]{eso-pic}
\usepackage{calc}
\usepackage{bookmark}
\usepackage{lastpage}
\usepackage{changepage}
\usepackage{paracol}
\usepackage{ifthen}
\usepackage{needspace}
\usepackage{iftex}

\ifPDFTeX
    \input{glyphtounicode}
    \pdfgentounicode=1
    \usepackage[T1]{fontenc}
    \usepackage[utf8]{inputenc}
    \usepackage{lmodern}
\fi

\usepackage{XCharter}

% Some settings:
\raggedright
\AtBeginEnvironment{adjustwidth}{\partopsep0pt}
\pagestyle{empty}
\setcounter{secnumdepth}{0}
\setlength{\parindent}{0pt}
\setlength{\topskip}{0pt}
\setlength{\columnsep}{0.15cm}
\pagenumbering{gobble}

% ── Section title ──
\titleformat{\section}{\needspace{2\baselineskip}\bfseries\large}{}{0pt}{}[\vspace{1pt}\titlerule]
\titlespacing{\section}{-1pt}{0.15 cm}{0.08 cm}

\renewcommand\labelitemi{$\vcenter{\hbox{\small$\bullet$}}$}

% ── Bullet lists ──
\newenvironment{highlights}{
    \begin{itemize}[
        topsep=0.10 cm,
        parsep=0.06 cm,
        partopsep=0pt,
        itemsep=0.06 cm,
        leftmargin=12pt
    ]
}{
    \end{itemize}
}

\newenvironment{highlightsforbulletentries}{
    \begin{itemize}[
        topsep=0.10 cm,
        parsep=0.06 cm,
        partopsep=0pt,
        itemsep=0.04 cm,
        leftmargin=12pt
    ]
}{
    \end{itemize}
}

\newenvironment{onecolentry}{
    \begin{adjustwidth}{0 cm + 0.00001 cm}{0 cm + 0.00001 cm}
}{
    \end{adjustwidth}
}

\newenvironment{twocolentry}[2][]{
    \onecolentry
    \def\secondColumn{#2}
    \setcolumnwidth{\fill, 4.5 cm}
    \begin{paracol}{2}
}{
    \switchcolumn \raggedleft \secondColumn
    \end{paracol}
    \endonecolentry
}

\newenvironment{threecolentry}[3][]{
    \onecolentry
    \def\thirdColumn{#3}
    \setcolumnwidth{, \fill, 4.5 cm}
    \begin{paracol}{3}
    {\raggedright #2} \switchcolumn
}{
    \switchcolumn \raggedleft \thirdColumn
    \end{paracol}
    \endonecolentry
}

\newenvironment{header}{
    \setlength{\topsep}{0pt}\par\kern\topsep\centering\linespread{0.6}
}{
    \par\kern\topsep
}

\let\hrefWithoutArrow\href

\begin{document}
    \newcommand{\AND}{\unskip
        \cleaders\copy\ANDbox\hskip\wd\ANDbox
        \ignorespaces
    }
    \newsavebox\ANDbox
    \sbox\ANDbox{$|$}

    \begin{header}
        \fontsize{16 pt}{24 pt}\selectfont Tanmay Bhuskute

        \vspace{3 pt}

        \normalsize
        \mbox{ +1 (732) 858-3027 }%
        \kern 1.0 pt%
        \AND%
        \kern 1.0 pt%
        \mbox{ tanmay.v.bhuskute@gmail.com }%
        \kern 1.0 pt%
        \AND%
        \kern 1.0 pt%
        \mbox{ github.com/Tanmayb05 }%
        \kern 1.0 pt%
        \AND%
        \kern 1.0 pt%
        \mbox{ linkedin.com/in/tanmay-bhuskute }%
        \kern 1.0 pt%
        \AND%
        \kern 1.0 pt%
        \mbox{ Cincinnati, OH }%
    \end{header}

    \vspace{2 pt}

%═══════════════════════════════════════
\section{SUMMARY}
%═══════════════════════════════════════
\begin{onecolentry}
Software Engineer and M.S. Computer Science candidate (Dec 2026) with hands on object-oriented programming, automated testing, and quality engineering experience at Siemens. Skilled in C/C++, Python, and Go, building test frameworks, CI/CD quality gates, and validation tooling for production infrastructure.
\end{onecolentry}

%═══════════════════════════════════════
\section{EDUCATION}
%═══════════════════════════════════════
\begin{twocolentry}{
    Aug 2024 - Dec 2026
}
    \textbf{Master of Science in Computer Science} \\
    University of Cincinnati
\end{twocolentry}
\begin{onecolentry}
Relevant Coursework: Advanced Algorithms, Database Theory, Object-Oriented Design.
\end{onecolentry}
\vspace{0.15 cm}

\begin{twocolentry}{
    Aug 2018 - Jul 2022
}
    \textbf{Bachelor of Engineering in Computer Engineering} \\
    AISSMS College of Engineering
\end{twocolentry}
\begin{onecolentry}
Relevant Coursework: Object-Oriented Programming, Data Analytics, Machine Learning.
\end{onecolentry}
%═══════════════════════════════════════
\section{TECHNICAL SKILLS}
%═══════════════════════════════════════
\begin{onecolentry}
\textbf{Programming:} C/C++, Java, Python, Go, JavaScript, TypeScript, SQL.
\end{onecolentry}
\vspace{0.06 cm}
\begin{onecolentry}
\textbf{Testing:} Automated Test Frameworks, Unit/Integration Testing, Terratest, JUnit, Pytest, Object-Oriented Design, Test-Driven Development.
\end{onecolentry}
\vspace{0.06 cm}
\begin{onecolentry}
\textbf{Systems \& Infrastructure:} Linux, Embedded/Runtime Environments, REST APIs, Microservices, Agile/Scrum, SDLC.
\end{onecolentry}
\vspace{0.06 cm}
\begin{onecolentry}
\textbf{Developer Tools:} Git, GitLab CI/CD, Jenkins, AWS, Docker, Kubernetes, Terraform, Ansible, SonarQube, Bash.
\end{onecolentry}
\vspace{0.06 cm}
\begin{onecolentry}
\textbf{Quality \& Collaboration:} Code Review, CI/CD Quality Gates, Cross-team Communication, Documentation, Mentoring.
\end{onecolentry}

%═══════════════════════════════════════
\section{WORK EXPERIENCE}
%═══════════════════════════════════════
\begin{twocolentry}{Jul 2022 -- Jun 2024}
\textbf{Software Development Engineer} ---- Siemens Industry Software
\end{twocolentry}
\begin{highlights}
  \item Served as Quality Coach and SonarQube Specialist, writing automated tests and enforcing CI/CD quality gates that raised automated code coverage from \textbf{35\% to 78\%} across the Agile Release Train.
  \item Designed \textbf{Ansible-based Blue-Green deployment} test and rollback framework, verifying \textbf{44 production environment upgrades} in 4 months across Linux, EC2, and Teamcenter X services with zero unresolved defects.
  \item Automated Teamcenter X infrastructure workflows with Python (Boto3) using object-oriented design, replacing legacy Ruby scripts and reducing deployment time by \textbf{40\%} across AWS environments.
  \item Migrated \textbf{40+ production environments} to Kubernetes using ArgoCD, Helm, and GitLab CI/CD, building repeatable test and validation pipelines for each release.
  \item Mentored \textbf{5 software engineers} on testing standards and structured onboarding, improving attention to detail and code quality across the team.
\end{highlights}

\vspace{0.15 cm}

\begin{twocolentry}{Mar 2022 -- Jun 2022}
\textbf{Software Developer Intern} ---- Siemens Industry Software
\end{twocolentry}
\begin{highlights}
  \item Developed \textbf{Go + Terratest} automated test framework for infrastructure code, catching \textbf{95\%} of misconfigurations pre-production before they reached runtime environments.
  \item Automated provisioning and validation of AWS resources with \textbf{Terraform}, reducing deployment time by \textbf{30\%} while ensuring 100\% successful test runs.
\end{highlights}

\vspace{0.15 cm}

\begin{twocolentry}{Sep 2021 -- Dec 2021}
\textbf{Software Developer Intern} ---- Government of Maharashtra, Water Resource Department
\end{twocolentry}
\begin{highlights}
  \item Built and deployed a full-stack \textbf{React and MySQL} workflow application using object-oriented design, digitizing departmental record management and improving operational efficiency by \textbf{40\%}.
\end{highlights}

%═══════════════════════════════════════
\section{PROJECTS AND RESEARCH}
%═══════════════════════════════════════
\begin{twocolentry}{Feb 2026 -- Present}
\textbf{Android Bug Reproduction Research} [Python, ADB, uiautomator2, scikit-image]
\end{twocolentry}
\begin{highlights}
  \item Built a device automation test harness (uiautomator2/ADB) with stall detection for a two-stage LLM pipeline reproducing Android GUI bugs, achieving \textbf{25.7x latency reduction} over the ViBR baseline across 53 runs on 20+ apps.
\end{highlights}

\vspace{0.12 cm}

\begin{twocolentry}{May 2026}
\textbf{PolicyPilot -- Medical Benefit Policy Intelligence Portal} [React, TS, Node, Zod, Vitest]
\end{twocolentry}
\begin{highlights}
  \item Designed 17-tool MCP server with Zod schema validation at every data boundary and Vitest unit/integration test coverage, eliminating a class of runtime type-mismatch bugs across all external tool calls.
\end{highlights}

%═══════════════════════════════════════
\section{EXTRACURRICULAR ACTIVITIES}
%═══════════════════════════════════════
\begin{onecolentry}
\begin{highlights}
  \item \textbf{Finalist, RevolutionUC 2025 Hackathon} (300 participants, 31 schools) \,|\, \textbf{SAFe Practitioner 5.1} (Mar 2023).
\end{highlights}
\end{onecolentry}

\end{document}
